% ============================================================================
% 5. Discussion  |  6. Limitations  |  7. Conclusion
% ============================================================================
\section{Discussion}\label{sec:discussion}

\subsection{Principal interpretation}\label{sec:interpretation}
The central finding is a \emph{pattern of conditionality}, now mapped along
three axes. Along the task axis, the standard $\sqrt{\dk}$ divisor is
significantly better than both alternatives on AG News at a 2-epoch budget~---
with large and medium paired effect sizes~--- yet statistically
indistinguishable from them on two binary sentiment datasets. Along the budget
axis, its advantage on binary tasks dissolves by 25 epochs, and the Stage-2
replica completion shows that the AG News effect itself does not survive a
fivefold reduction in data and sequence length: under matched seeds the entire
replica grid (core, $\dk$ grid, convergence training, corrected language
modelling) is variant-invariant to within measurement resolution. Along the
operating-point axis~--- the axis Stage 3 adds~--- the nulls themselves are
regime facts: a quadratic power law on the raw logit scale predicts that the
divisor contrast is inert below $\sigma \approx 1$ and decisive above it, the
replica grids sit at $\sigma \approx 0.3$--$1.1$ (measured KL $\le 0.006$
nats), forcing $\sigma$ an order of magnitude higher flips the sign of the
effect at $g = 32$, and a pretrained GPT-2 sitting at $\sigma \approx 45$--63
loses a factor of 94 in perplexity when the same constant is swapped to
$1/\dk$ at inference. A one-line architectural constant inherited from
\cite{vaswani2017} is therefore best understood not as a universally optimal
choice but as a default whose benefit is confined to a specific operating
point~--- full data scale, early in training, on multi-class topic
classification~--- and whose influence vanishes once the softmax operates
near uniformity, while at frontier operating points it is the alternatives
that become untenable.

\subsection{Scale dependence and the near-uniform regime}\label{sec:scale}
The Stage-2 null results are not a failure of the manipulation: the
diagnostics show the mechanism is present but quantitatively muted. At replica
scale the raw pre-division logit standard deviation is of order 1, so after
dividing by $\sqrt{\dk}$ (5.7--11.3) or $\dk$ (32--128) every variant's logits
sit well inside the softmax's linear (near-uniform) regime, attention entropy
is pinned at 4.05 of a possible 4.16 nats for all variants, and the Q/K
gradient pathway carries two orders of magnitude less signal than the value
pathway~--- an imbalance that AdamW's per-parameter scale normalisation then
equalises away. Two consequences follow. First, the variance argument for
$\sqrt{\dk}$ is a statement about logit variance \emph{relative to the
softmax's dynamic range}; it predicts differences only in regimes where the
unscaled logit variance is comparable to the divisor, which small-data,
short-budget runs do not reach. Second, claims about scaling functions are
claims about an operating point, and should be reported together with the
logit scale that makes the divisor bite.
Stage 3 turns this diagnosis from an explanation into a measurement. The
power law of Eq.~\eqref{eq:law} quantifies the bite region; the gain sweep
shows the outcome metric responding on cue once the operating point enters it
(flat to $g \le 16$, $-1.4$\,pp for the standard divisor at $g = 32$); the
regime-restored runs show that at $d_k = 32$ even $12\times$ the step budget
leaves $\sigma$ at 1.1 (inertia is real), while at $d_k = 128$ the larger
head drives $\sigma$ to $\approx 5$ and the Stage-1 ordering begins to
re-emerge; and the pretrained decoder probe shows the same constant being
\emph{decisive} ($+27$ to $+9{,}253\%$ perplexity) at GPT-2's native
operating point. The Stage-1 AG News effect, observed precisely where raw
logit variance comes closest to the divisor, remains plausible under this
account~--- but until it is reproduced at full scale it should be treated as
provisional rather than established.

\subsection{Why would AG News care about the divisor?}\label{sec:why-agnews}
Three observations combine into a coherent account. First, the variance argument
behind $\sqrt{\dk}$ \cite{vaswani2017} implies that the $1/\dk$ divisor shrinks
the logit standard deviation by a further factor of $\sqrt{\dk} \approx 5.7$;
at initialization the attention logits of the $1/\dk$ variant have a standard
deviation of order $1/\sqrt{\dk} \approx 0.18$ rather than 1, pushing the softmax
towards uniformity. Second, the measured diagnostics match this account:
attention entropy rises and maximum attention weight falls monotonically with
$s$ (Section~\ref{sec:diagnostics}). Third, AG News is the task where
failure-to-differentiate should bite hardest: it is 4-class, short-document, and
topic-level, so early error correction plausibly benefits from a few sharp,
content-word attention links, whereas binary sentiment over long reviews can be
solved with diffuse pooling-like attention. The inter-head cosine similarity
trend (0.23 $\to$ 0.42 as $s$ grows) adds that flatter scaling also erodes
head diversity \cite{voita2019,michel2019}, reducing the ensemble capacity that
multi-class discrimination can draw on. The mixed-scale ablation (E1), where
keeping two sharp heads helped on SST-2, hints that diversity per se~--- not the
global constant~--- is the operative resource.

\subsection{Sharper attention as a free improvement direction}\label{sec:sharper}
The temperature ablation points in the opposite direction: on both sentiment
datasets the sharpest tested temperature ($\tau = 0.5$, i.e.\ doubling the logit
scale relative to standard) produced the highest means, by about $+1.6$\,\pp{} on
SST-2. This aligns with the distillation literature on attention temperature
\cite{zhang2022temp,hinton2015} and with Scalable-Softmax's motivation
\cite{nakanishi2025} that the effective temperature should not be treated as a
universal constant. Given that temperature is parameter-free and computationally
free (Section~\ref{sec:efficiency}), and that the 2-epoch AG News result shows
temperature effects \emph{can} reach significance under matched conditions, the
practical recommendation is straightforward: when training small transformers on
a new task, sweep $\tau$ around 1 (e.g.\ $\{0.5, 1, 2\}$) under a seed-matched
design before freezing the architecture. The cost is negligible; the present
data suggest the yield is real, if task-dependent.

\subsection{The early-training effect and its decay}\label{sec:early}
On the binary tasks, extended training erased the differences (Section
\ref{sec:budget}), and the Stage-2 convergence runs sharpen the picture: with
model selection by validation loss, every variant's validation loss bottoms at
the same epoch ($\approx$ 2 on SST-2, $\approx$ 1.6 on IMDB), early stopping
halts at epochs 5--6, and the learning curves are indistinguishable at every
epoch. Combined with the entropy-collapse literature
\cite{zhai2023,dong2021}, a natural reading is that the divisor matters most
while the attention distribution is still being shaped; once representations
mature, the network can compensate for the initial sharpness of its attention.
The learning-rate compensation ablations (E2/E4) add a useful negative result:
the $1/\dk$ deficit is \emph{not} equivalent to a smaller effective learning rate
that a per-block LR rescaling can restore~--- and the multi-seed replica
ablations confirm the single-seed pattern (LR compensation neutral-to-harmful,
mixed-scale heads neutral), while failing to reproduce the Stage-1 E1 excursion
of $+1.6$\,\pp{} for mixed-scale heads on SST-2.

\subsection{Implications for experimental practice}\label{sec:practice}
Methodologically, this study is a demonstration that cheap architectural
comparisons deserve the seed-matched treatment. Between-run comparisons at 2
epochs on these datasets would need to overcome standard deviations of about
1\,\pp{} (SST-2, IMDB) with between-run variance; the paired design recovers
resolvable effects from 25 seeds even where the raw means differ by only
0.1--0.3\,\pp{} (AG News). Given the prevalence of single-seed architectural
claims in the scaling literature \cite{bouthillier2021,dodge2019,dror2018}, the
protocol used here~--- bit-identical initialization, shared data order, paired
tests, Holm correction, effect sizes~--- is directly reusable by practitioners
evaluating any parameter-free modification (temperature, biases, clamps,
positional offsets) and costs no additional compute beyond a handful of extra
seeds.

\subsection{Threats to validity}\label{sec:threats}
Beyond the enumerated limitations, three threats deserve note. \emph{Construct}:
test accuracy on subsampled training sets may not capture the behaviours that
matter at scale; we mitigate by also measuring perplexity and mechanistic
diagnostics, and by reporting both a full-scale and a replica stage.
\emph{Internal}: all variants share code paths except one division~--- the
manipulation is auditable in a single line~--- and unit tests plus the
Stage-2 mask-verification suite pin down initialization identity and exact-zero
masked attention, so the internal threat is minimal. The Stage-2 null results
are consistent across five independent grids, which argues against a wiring
error: a variant applied to the wrong tensor would not produce identical
entropy \emph{and} identical accuracies \emph{and} the Stage-1 gradient-ratio
signatures. \emph{External}: results at $d_{\text{model}} = 128$, 2 layers, and
64--128-token sequences may not extrapolate to large models; the Stage-3
programme addresses this directly~--- the power law makes the extrapolation
quantitative, the pretrained-decoder probe confirms its direction and
magnitude at 124M--355M scale, and the shipped GPU script extends the identical
protocol to 1B--7B, but the trained (not inference-only) behaviour of models
in that range remains to be measured.

% ============================================================================
\section{Limitations}\label{sec:limitations}
The study has five principal limitations.
\begin{enumerate}[leftmargin=1.9em]
  \item \textbf{Replica scale.} Stage 2 documents a deliberate fivefold
        reduction in data (2{,}000 vs.\ 10{,}000 examples) and sequence length
        (64 vs.\ 128), with a 96k-token LM corpus. The replica null results are
        statements about that operating point; they bound the Stage-1 AG News
        effect's generality but cannot confirm or refute it at full scale.
  \item \textbf{Model scale.} A 2-layer, $d_{\text{model}} = 128$ encoder is
        three orders of magnitude smaller than frontier models. Stage 3
        narrows this gap as far as the CPU sandbox allows: the operating-point
        law is validated in simulation, the gain sweep and regime-restored
        runs traverse the bite region from below, frontier stabilisers are
        tested in context, and pretrained GPT-2 (124M, 355M) is probed at its
        native operating point, where the divisor contrast is enormous.
        What remains untested is \emph{trained} behaviour at 1B--7B: the
        shipped GPU script covers the inference-time intervention and
        diagnostics, but fine-tuning or training at those scales, and the
        interaction of the divisor with QK-Norm or soft-capping \emph{during
        training}, are beyond this study's compute.
  \item \textbf{LM budget.} The corrected LM grid uses 96k training tokens and
        3 epochs (perplexity $\approx$ 873); the leak probes show the v1
        artifact's direction, but its fully collapsed state (ppl $\approx$ 1.08
        at 25{,}000 GPU steps) is not reproduced at CPU scale. The Stage-3 LM
        gain arm (sequence length 128, $g \le 4$) likewise operates in the
        inert region, so the LM-scale divisor question remains open.
  \item \textbf{Task scope.} Four text datasets; vision, speech, code, and
        long-context regimes are untouched, as are encoder--decoder and
        decoder-only architectures.
  \item \textbf{Full-scale gaps.} The Stage-1 extended budget (30 of 450
        planned full-scale runs) and the full-scale E1--E4 multi-seed arms
        remain incomplete; the sequence-length axis of the sensitivity grid is
        covered only implicitly (64-token replica vs.\ 128-token Stage 1).
        These are the principal targets for the next full-scale campaign.
\end{enumerate}

% ============================================================================
\section{Conclusion}\label{sec:conclusion}
Under a seed-matched paired design with 787 completed runs across a GPU-scale
stage, a replica-scale completion, and an operating-point programme, this
study measured what one line of the Transformer~--- the $\sqrt{\dk}$ divisor
in scaled dot-product attention~--- actually contributes. The answer is
sharply conditional in a now-quantitative sense. At full data scale and a
2-epoch budget, the standard divisor outperformed $1/\dk$ and temperature
$\tau = 2$ on 4-class AG News ($+0.31$ and $+0.10$\,\pp, $p_{\text{Holm}} <
0.05$, $d_z$ up to 1.52), but was indistinguishable on binary sentiment.
Elsewhere the scaling function is empirically inert: the replica completion of
the core grid, the $\dk \in \{32,64,128\}$ grid, train-to-convergence with
early stopping, and the corrected causal-mask language-modelling grid (25
seeds per arm, perple\-xity $\approx$ 873 with a verified exact-zero mask) all
show differences within noise, and the original near-1 perplexities are
diagnosed and reproduced-in-direction as a mask-leak artifact. Stage 3
explains the inertness and bounds its reach. A second-order power law
predicts that two divisors induce attention distributions separated by
$\mathrm{KL} \approx \tfrac{1}{2}\Delta\beta^2\sigma^2(1-1/n)$: the replica
grids sit at measured KL $\le 0.006$ nats, $12\times$ longer training leaves
the operating point at $\sigma = 1.1$, injecting logit variance until
KL $= 0.64$ nats makes the 2017 default the \emph{worst} variant
($-1.4$\,\pp, $d_z = -1.76$), frontier stabilisers bound the logits but not
the contrast, and in pretrained GPT-2 ($\sigma \approx 45$--63) the same
no-retraining swap moves perplexity by $+27$--$55\%$ ($\tau = 2$) to
$40$--$93\times$ ($1/\dk$). The practical conclusions are that $\sqrt{\dk}$
remains a sound default, that the softmax temperature is a cheap knob worth
tuning per task, that masking bugs in language-modelling pipelines are best
caught by physically-motivated sanity floors on perplexity, that any claim
about a parameter-free attention modification should report the operating
point it was tested at (one forward pass suffices to measure $\sigma$), and
that matched-seed paired designs are the cheap instrument that makes such
claims testable.

\section*{Data and code availability}
The complete reproducibility package~--- source code, YAML configurations, unit
tests, per-seed JSON records, statistical analysis outputs, figure-generation
scripts, and this manuscript~--- is publicly available under an MIT licence at
\url{https://github.com/razahasan2000/Attention-Scaling}. Supplementary material
with full hyperparameters, per-seed summaries, diagnostics, and reproduction
instructions accompanies this paper.

\section*{Declaration of interests}
The author declares no competing interests.
